% Delegated research by GPT-6 Astra (OpenAI), 20 November 2026.
% This is an INCOMPLETE contribution. Statuses are explicit in results.csv.
\section{Seshadri, henselian comparison, and the structure of the Picard functor}
\label{sec:bce}

\subsection{The distinction between the manuscript and the supplied claims}

At 57:110 the question retains the hypothesis of essential properness of the
\emph{Picard functor}; replacing the representable Picard scheme by its functor
does not discard that properness hypothesis. The supplied Part B1 and Lead H2
discard it. The counterexample below disproves this broader formulation.
At 57:111 the text is unambiguous:
\begin{quote}
``une intersection compl\`ete relative de dim relative $\geq3$''.
\end{quote}
The complete-intersection hypothesis must likewise be retained in the positive
$R_2$ criterion. Finally, 57:314 says
\begin{quote}
``Prouver, si $X\to S$ est normal, que $\underline{\mathrm{Pic}}^0_{X/S}$ est
`ferm\'e' dans $\underline{\mathrm{Pic}}_{X/S}$ mieux, que
$\underline{\mathrm{Pic}}^0_{X/S}$ est propre sur $S$ ??''
\end{quote}
Here normality is a property of the \emph{morphism}: flatness with geometrically
normal fibres. The shorter wording on 57:315 must not silently be substituted
by mere normality of the total space.

\subsection{B1: a proper counterexample in every characteristic}

Let $k$ be algebraically closed, let $R=k[[t]]$, and put
\[
 X=\operatorname{Proj}R[x,y,z,w]/(xy-z(z-tw)),\qquad D=V(x,z).
\]
The equation is primitive over the DVR, and its quotient ring is torsion-free;
thus $X/R$ is flat. It is projective. For $t\ne0$, the change of homogeneous
coordinate $v=z-tw$ identifies its fibre with the smooth quadric $xy=zv$.
The special fibre is the integral quadric cone $xy=z^2$ in $\mathbf P^3$.
It is a hypersurface and hence Cohen--Macaulay, and its only singular point is
the vertex $[0:0:0:1]$, also in characteristic two. It is therefore normal and
$S_2$. These assertions persist after extending $k$.

The divisor $D$ is $\operatorname{Proj}R[y,w]\simeq\mathbf P^1_R$, so it has no
vertical component, and its generic fibre is Cartier on the smooth quadric.
On the chart $w=1$, at the origin, its prime ideal is $(x,z)$ in
\[
 B=\bigl(R[x,y,z]/(xy-z^2+tz)\bigr)_{(t,x,y,z)}.
\]
The ring $B$ is normal: it is a hypersurface with an isolated singularity.
All its defining relations have order at least two, whereas $x,z$ have
linearly independent images in $\mathfrak m/\mathfrak m^2$. Since
$\mathfrak m(x,z)\subset\mathfrak m^2$, the ideal $(x,z)$ needs two generators.
It is not invertible in this local ring. Consequently $D$ is not Cartier.
Taking the product with $\mathbf P^1_R$ also gives relative dimension three.

The sign in the supplied R1 example is correct when $\operatorname{char}k\ne2$:
if $i^2=-1$, then $(z-it)(z+it)=z^2+t^2$, so $(x,z-it)$ is indeed a prime
divisor on $xy-z^2-t^2=0$. The split projective version above removes the
restriction on the characteristic.

\subsection{B2: sufficiency, nonnecessity, and the exact obstruction}

Suppose $S$ is regular and integral, $f:X\to S$ is flat of finite type and a
relative local complete intersection, and its geometric fibres satisfy $R_2$.
The fibres are Cohen--Macaulay, hence normal, and $X$ is a normal absolute
local complete intersection. Let $D$ be a Weil divisor Cartier on $X_\eta$.
At any point of $X$ of codimension at most three lying over a nongeneric point
of $S$, the fibre-local dimension is at most two. The fibre is regular there,
so flatness and regularity of the base make the local ring of $X$ regular.
At points over $\eta$, $D$ is Cartier by hypothesis. Thus the non-Cartier locus
of $D$ has codimension at least four.

If it is nonempty, choose a generic point $x$ of that locus. On
$\operatorname{Spec}\mathcal O_{X,x}$ the divisor is Cartier away from the
closed point. This local ring is a complete intersection of dimension at least
four. Parafactoriality (SGA~2, XI, Theorem~3.13(ii)) makes the corresponding
line bundle on its punctured spectrum trivial. Normality then extends a
principal Weil divisor from the punctured spectrum, forcing $D$ to be Cartier
at $x$, a contradiction.

The $R_2$ condition is not necessary, even within the proper scope of B1.
Over $R=k[[t]]$ in characteristic different from two, compare the projective
families in $\mathbf P^3_R$
\[
 X_1: xy-z^2-tw^2=0,\qquad X_2: xy-z^2-t^2w^2=0.
\]
The only singular point of either special fibre is its cone vertex $w=1$.
At that point the derivative of the first equation with respect to $t$ is
$-1$, so the total space $X_1$ is regular everywhere and locally factorial.
Its special fibre is a normal surface node, not $R_2$, but every divisor on
$X_1$ is Cartier. The family $X_2$ has the non-Cartier horizontal divisor
$(x,z-itw)$, as in R1. Both special fibres are the same quadric cone, and
both geometric generic fibres are smooth projective quadrics. Taking the
product of both families with $\mathbf P^1_R$ gives relative dimension three
and preserves all these properties. Thus even proper relative complete
intersections with the same base and abstract geometric fibre types can
have different answers: those data alone do not determine the Cartier
extension property.

Without the relative complete-intersection assumption, even $R_2$ is not
sufficient. In $\mathbf P^6_R$ with coordinates $[a:b:c:d:e:f:w]$ impose
\[
 \operatorname{rank}\begin{pmatrix}a&b&c\\d&e&f\end{pmatrix}\leq1,
 \qquad a+e=tw.
\]
The generic fibre is the Segre variety $\mathbf P^1\times\mathbf P^2$.
The special fibre is the projective cone over the smooth hyperplane section
$u_0v_0+u_1v_1=0$ of that Segre variety. It is a normal Cohen--Macaulay
threefold with just its vertex singular, hence $R_2$. Flatness follows from
the constant Hilbert polynomial of the hyperplane sections of the projective
cone over the Segre variety (equivalently from the regular sequence defined
by $t$ and the hyperplane equation). At the vertex $w=1$, the completed local
ring of the total space, after eliminating $t=a+e$, is the rank-one
$2\times3$ determinantal cone. This ring is normal. The divisor
$D=V(a,b,c)$ is $\mathbf P^2_R$, with $e=tw$, and is generically Cartier.
Its ideal at the vertex requires three generators, since the determinantal
relations are quadratic. It is not Cartier. The same construction with a
$2\times(m+1)$ matrix, $m\geq2$, has fibres of dimension $m+1$ whose isolated
singularity satisfies $R_m$. No fixed finite Serre regularity condition alone
replaces the complete-intersection hypothesis.

For a fixed family there is an exact, deformation-sensitive criterion. Suppose
$S$ is regular integral and $X/S$ is flat with geometrically integral normal
fibres. Every vertical prime divisor on $X$ is the pullback, with multiplicity
one, of a prime divisor of $S$, and hence is Cartier. Therefore every horizontal
Weil divisor whose generic restriction is Cartier is Cartier if and only if
\[
 \ker\left(\operatorname{Cl}(X)/\operatorname{Pic}(X)
 \longrightarrow \operatorname{Cl}(X_\eta)/\operatorname{Pic}(X_\eta)\right)=0.
\tag{B.1}
\]
Indeed split any Weil divisor into its horizontal and vertical parts; the
vertical part does not change its Cartier obstruction. If $X_\eta$ is locally
factorial, (B.1) is simply local factoriality of $X$. This is an exact criterion
for the given family, not an assertion that the requested classification by
abstract fibres has been produced.

